\section{Approach}
\label{sec:approach}

\begin{figure*}[t]
    \centering
    \includegraphics[width=\textwidth]{figs/framework.pdf}
    \caption{Illustration of ZSCL in feature space. Fig. 3(a) shows the pipeline of distillation. The original and the current model encode the reference images and texts, respectively. The probability distribution of images and texts with respect to each other is distilled. Fig. 3(b) displays how distillation loss preserves the feature space. Compared with the reference dataset, the features of fine-tuning tasks lie in a small subspace. The distillation loss preserves the structure of the feature space by maintaining relative distances.}
    \label{fig:framework}
\end{figure*}


\subsection{Preliminaries}

\paragraph{Continual Learning.} Given $n$ tasks $[\mathcal{T}^1, \mathcal{T}^2, \cdots, \mathcal{T}^n]$, continual training is conducted in sequence on each task $\mathcal{T}^i=(\mathcal{D}^i, C^i), i=1,\dots,n$. Here, $\mathcal{D}^i$ represents the task dataset $\{(\bm{x}^i_j, \bm{y}^i_j\}_{j=1}^{N_i}$, where $\bm{x}^i_j$ is an image, $\bm{y}^i_j$ is a one-hot vector indicating the ground truth, and $N_i$ is the number of images in the dataset. Class names $C^i=\{c^i_j\}_{j=1}^{m_i}$ maps the label of an image to an object name, where $m_i$ is the number of classes for task $\mathcal{T}^i$. The objective of continual learning is to achieve good performance on all tasks. 

Two continual learning settings are focused on in this study~\cite{van2019three}. In task-incremental learning, at inference, the image $\bm{x}$ to be predicted is given with its task identity $t$, so the model only needs to distinguish between different classes in $C^t$. In class-incremental learning, the task identity $t$ is not given. Thus we need to predict with the combined class set $C=\bigcup_{i=1}^nC^i$.

\paragraph{CLIP model.} The CLIP model contains an image encoder $f_i$ and a text encoder $f_t$. The inference process of the CLIP model for image classification is as follows. First, for task $\mathcal{T}^i$, each class $c$ is transformed into a sentence by a template like ``a photo of \{$c$\}''. Then $f_t$ encodes the classes into text embeddings $\{\bm{t}^i_j\}_{j=1}^{m_i}$. An image encoder encodes input image $\bm{x}_k$. The similarity score between the image embedding and text embeddings are calculated as $\bm{s}^i_{k,j}=\Braket{f_i(\bm{x}_k),\bm{t}^i_j}$, where $\Braket{\cdot,\cdot}$ denotes the cosine similarity. The class with the largest similarity score is the prediction for the image. 

To fine-tune the CLIP model for downstream tasks, cross-entropy loss $\text{CE}$ is applied to the similarity score with a temperature scaling:
\begin{equation}
    \mathcal{L}_{\text{CE}} = \frac{1}{N}\sum_{i=1}^N\text{CE}(\tau\cdot \bm{s}_i, \bm{y}_i),
\end{equation}
where $\tau$ is a parameter learned during the pre-training.

\subsection{Distillation in Feature Space}

Well-learned feature space for aligned images and texts enables vision-language models' strong zero-shot transfer ability. It also facilitates the learning of downstream tasks. Compared with the pre-training dataset, downstream datasets lie in a small scope in the feature space (shown in \cref{fig:framework}(b)). Direct fine-tuning of downstream tasks greatly distorts the feature distribution of out-of-distribution data, which leads to a significant drop in zero-shot prediction performance.

While the cross-entropy improves the performance by altering fine-tuned feature subspace, we need a new regularization to preserve the potential out-of-distribution feature space. The relative similarity between one image and different texts is:
\begin{equation}
    \bm{p}=\text{Softmax}(\bm{s}_1, \cdots, \bm{s}_{m}).
\end{equation}
We hope the above similarity distribution is stable during fine-tuning for all potential images and texts. Given a teacher model $\overline{f}$, distillation loss can be applied to penalize changes from the original distribution:
\begin{equation}
    \mathcal{L}_{\text{dist\_img}}=\text{CE}(\bm{p},\bm{\overline{p}})=-\sum_{j=1}^m\bm{p}_j\cdot\log\bm{\overline{p}}_j, \label{eq:distimg}
\end{equation}
where $\overline{\bm{p}}$ is the distribution given by the teacher model.

Although the distillation form has been widely used in previous methods~\cite{hinton2015distilling,li2017learning,rebuffi2017icarl}, they are applied to continual learning from scratch. We investigate different components of the distillation loss for enhancing pre-trained vision-language models.

Three components are discussed in this paper in detail: the data source, the teacher model, and the loss design. \cref{subsec:main_properties} shows the performance for the different choices. First, for the data source, LwF~\cite{li2017learning} uses data from the current task, while iCarl~\cite{rebuffi2017icarl} carefully selects data from previous tasks. However, data from downstream tasks span a small subspace and are not good enough to preserve the whole feature space. An ideal choice is the pre-training dataset. However, the CLIP pre-training dataset is private, and the size is too large to load. To solve this challenge, we introduce the reference dataset. A reference dataset is a publicly available image dataset with enough semantics. Enough semantics can be seen as random sampling in the whole feature space. The texts can be related class names, unrelated sentences, or even random tokens. This is because text semantics are easier to sample, and we need not ground truth to calculate the distance between one image to sufficient text samples spread among feature space.

For the teacher model, \cite{li2017learning,rebuffi2017icarl} use the model after learning task i-1 and before learning task i as the teacher model. During the continual training of the vision-language model, the feature space deviates gradually from the initial model. Using fine-tuned models as teacher models enlarges this change. In contrast, we find using the pre-trained model as a teacher model not only preserves the zero-shot transfer ability but also takes advantage of well-learned feature space for better downstream performance.

Finally, previous distillation loss is applied on traditional backbones, where a classification head gives the probability for different labels. For the vision-language model, the probability is calculated based on the relative distance between images and texts. Thus, in addition to \cref{eq:distimg}, we impose regularization $\mathcal{L}_{\text{dist\_txt}}$ on the distances from a text to a batch of images. The whole framework is shown in \cref{fig:framework} (a) with the following training loss: \begin{equation}
    \mathcal{L} = \mathcal{L}_{\text{ce}}+\lambda\cdot (\mathcal{L}_{\text{lwf\_img}} + \mathcal{L}_{\text{lwf\_txt}}).
\end{equation}

\subsection{Weight Ensemble in Parameter Space}

\begin{figure}[t]
    \centering
    \includegraphics[width=0.8\columnwidth]{figs/weightavg.pdf}
    \caption{Models during training contain different tradeoffs between zero-shot and new task performance. Points for FT are sampled every 100 iterations, and the ones for WiSE-FT means different $\alpha$ choices. WE ensembles models during training and achieve better performance.}
    \label{fig:weightavg}
\end{figure}


Machine learning models integrate learned knowledge in their parameters. To mitigate the forgetting problem, a series of works~\cite{kirkpatrick2017overcoming,zenke2017continual,aljundi2018memory} impose regularization losses on the changes of parameters. Weight consolidation (WC)~\cite{kirkpatrick2017overcoming} imposes the following loss: 
\begin{equation}
    \mathcal{L}_{\text{WC}}=\sum_i(\theta_i - \overline{\theta}_i)^2.
\end{equation}
where $\theta$ is the parameters of the current model, and $\overline{\theta}$ is the reference ones. Although this regularization prevents forgetting, it hinders learning new tasks in a challenging CL setting.

Apart from regularization losses, another method in parameter space is to ensemble different model weights. Model soup~\cite{wortsman2022model} averages weights of multiple fine-tuned models to improve the model's robustness but introduces additional training costs. WiSE-FT~\cite{wortsman2022robust} propose a weighted average between fine-tuned model and the original model to improve the out-of-distribution prediction robustness:
\begin{equation}
    f(x; (1-\alpha)\cdot \theta_0 + \alpha\cdot\theta_1),
\end{equation}
where $\theta_0$ is the original model and $\theta_1$ is the fine-tuned one. However, this method is hyper-parameter-sensitive where different $\alpha$ gives different tradeoffs between zero-shot transfer ability and downstream task performance (blue line in \cref{fig:weightavg}.

Inspired by this, we extend the weighted average to the CL setting. The motivation for the weighted average is to prevent fine-tuning from losing too much knowledge in the original model. As training goes by (green line in \cref{fig:weightavg}), the model performs better on new tasks while losing accuracy on previous ones. The models among training compose a sequence of different learning-forgetting tradeoffs. Instead of interpolating only between the initial and the final model, our method weight ensemble (WE) averages the weights in the sequence during the training time:
\begin{equation}
    \hat{\theta}_t = \begin{cases}
        \theta_0 & t=0 \\
        \frac{1}{t+1}{\theta}_{t} + \frac{t}{t+1}\cdot\hat{\theta}_{t-1} & \text{every I iterations}
    \end{cases}.
\end{equation}
where model weight sampling happens every $I$ iteration. 
The method is related to Stochastic Weight Averaging (SWA)~\cite{izmailov2018averaging}, but we do not use a modified learning rate schedule here because instead of getting better generalization ability, WE aim to give an improved learning-forgetting tradeoff. 
As shown in \cref{fig:weightavg}, WE achieves better performance on downstream tasks than WiSE-FT. In addition, while WiSE-FT is sensitive to different values of $\alpha$, our method is much more robust under different hyper-parameter ($I$) choices. 